% =====================================================================
\section{Energía}
\label{sec:energia}
% =====================================================================

\subsection{Segundo teorema de König}

\begin{teorema}[König para la energía cinética]
\label{teo:konig2}
\begin{equation}
  \resultado{T=\tfrac12\,MV^{2}+T',\qquad T'\defeq\sa\tfrac12\,m_{\alpha}{v'_{\alpha}}^{2}} .
  \label{eq:konig2}
\end{equation}
La energía cinética es la del CM, que lleva toda la masa, más la
\emph{interna} $T'$ del movimiento relativo al CM.
\end{teorema}

\begin{proof}
Con $\vb{v}_{\alpha}=\Vc+\vp_{\alpha}$, $v_{\alpha}^{2}=V^{2}+{v'_{\alpha}}^{2}+2\,\Vc\cdot\vp_{\alpha}$, y por tanto
\[
  T=\tfrac12V^{2}\sa m_{\alpha}+\sa\tfrac12 m_{\alpha}{v'_{\alpha}}^{2}
  +\Vc\cdot\underbrace{\sa m_{\alpha}\vp_{\alpha}}_{\vb{0}\ \eqref{eq:prop-cm}} .\qedhere
\]
\end{proof}

\begin{observacion}
$T'$ es la parte de la energía cinética que «no se ve» desde fuera: en un
gas es la agitación térmica. En un choque o una explosión, $\tfrac12MV^{2}$ no
cambia si $\vb{F}\ext=\vb{0}$, y toda la energía puesta en juego se reparte en $T'$.
\end{observacion}

\subsection{Teoremas de la energía}

\begin{lema}[Derivada de una distancia]
\label{lem:radial}
Sea $\vb{r}_{\alpha\beta}=r_{\alpha\beta}\,\uv{r}_{\alpha\beta}$, donde $\dot{\vb{r}}_{\alpha\beta}$ (vector) es la velocidad relativa
de $\alpha$ respecto de $\beta$ y $\dot r_{\alpha\beta}$ (escalar), el ritmo de cambio de su distancia. Entonces
\begin{equation}
  \resultado{\dot{\vb{r}}_{\alpha\beta}\cdot\uv{r}_{\alpha\beta}=\dot r_{\alpha\beta}} :
\end{equation}
la proyección de la velocidad relativa sobre la recta que une las partículas
es el ritmo al que cambia su distancia.
\end{lema}

\begin{proof}
\emph{Vía 1.} Derivando $r_{\alpha\beta}^{2}=\vb{r}_{\alpha\beta}\cdot\vb{r}_{\alpha\beta}$:
$2r_{\alpha\beta}\dot r_{\alpha\beta}=2\,\vb{r}_{\alpha\beta}\cdot\dot{\vb{r}}_{\alpha\beta}$, y se divide entre $2r_{\alpha\beta}$.

\noindent\emph{Vía 2 (base móvil).} Por la regla del producto,
\begin{equation}
  \dot{\vb{r}}_{\alpha\beta}=\underbrace{\dot r_{\alpha\beta}\,\uv{r}_{\alpha\beta}}_{\text{radial: cambia la distancia}}
  +\underbrace{r_{\alpha\beta}\,\dot{\uv{r}}_{\alpha\beta}}_{\text{perpendicular: giro}} ,
  \label{eq:radial-giro}
\end{equation}
y $\dot{\uv{r}}_{\alpha\beta}\perp\uv{r}_{\alpha\beta}$ porque $\uv{r}_{\alpha\beta}$ es unitario (Tema 1, lema del vector
unitario). Al multiplicar por $\uv{r}_{\alpha\beta}$, el término de giro se anula y queda $\dot r_{\alpha\beta}$.
\end{proof}

% ---------------------------------------------------------------------
\begin{figure}[ht]
\centering
\begin{tikzpicture}[scale=1.35]
  \coordinate (B) at (0,0);
  \coordinate (A) at (3.2,1.6);
  \pgfmathsetmacro{\ang}{atan2(1.6,3.2)}
  \pgfmathsetmacro{\rr}{sqrt(3.2^2+1.6^2)}
  % circunferencia r = const (sólo por debajo de alpha)
  \draw[texto!30,line width=0.45pt,dash pattern=on 1.5pt off 1.5pt]
    (B) ++({\ang-26}:\rr) arc[start angle={\ang-26},end angle={\ang-1.5},radius=\rr];
  % vector relativo
  \draw[texto!55,line width=0.6pt,-{Stealth[length=4pt,width=3pt]}] (B) -- (A)
    node[pos=0.5,anchor=north west,etiq,inner sep=2pt]{$\vb{r}_{\alpha\beta}$};
  % componentes
  \coordinate (Vr) at ($(A)+({\ang}:1.15)$);
  \coordinate (Vp) at ($(A)+({\ang+90}:0.95)$);
  \coordinate (V) at ($(Vr)+({\ang+90}:0.95)$);
  \draw[texto!25,line width=0.35pt,dash pattern=on 1.5pt off 1.5pt] (Vr) -- (V) -- (Vp);
  \draw[cian,line width=0.7pt,-{Stealth[length=4pt,width=3pt]}] (A) -- (Vr)
    node[anchor=north west,etiq,inner sep=1pt]{$\dot r_{\alpha\beta}\,\uv{r}_{\alpha\beta}$};
  \draw[rosa,line width=0.7pt,-{Stealth[length=4pt,width=3pt]}] (A) -- (Vp)
    node[anchor=south east,etiq,inner sep=1pt]{$r_{\alpha\beta}\,\dot{\uv{r}}_{\alpha\beta}$};
  \draw[ambar,line width=0.8pt,-{Stealth[length=4.5pt,width=3.2pt]}] (A) -- (V)
    node[anchor=south west,etiq,inner sep=1pt]{$\dot{\vb{r}}_{\alpha\beta}$};
  % ángulo recto mínimo
  \draw[texto!50,line width=0.3pt] ($(A)+({\ang}:0.11)$) -- ($(A)+({\ang}:0.11)+({\ang+90}:0.11)$) -- ($(A)+({\ang+90}:0.11)$);
  % partículas
  \fill[texto] (B) circle (1.4pt);
  \node[etiq,anchor=north east,inner sep=2pt] at (B) {$\beta$};
  \fill[texto] (A) circle (1.4pt);
  \node[etiq,anchor=north west,inner sep=3pt] at ($(A)+({\ang-90}:0.05)$) {$\alpha$};
  \node[peq,anchor=north,inner sep=2pt] at ($(B)+({\ang-26}:\rr)$) {$r_{\alpha\beta}=\text{const}$};
\end{tikzpicture}
\caption{Descomposición \eqref{eq:radial-giro} de la velocidad relativa: la
componente radial (azul) estira o acorta la distancia; la perpendicular
(azul claro) hace girar la pareja sobre la circunferencia $r_{\alpha\beta}=\text{const}$ (a
trazos) sin cambiar la distancia. Una fuerza central sólo «ve» la parte radial.}
\label{fig:radial-giro}
\end{figure}
% ---------------------------------------------------------------------

\begin{proposicion}[Variación de la energía cinética]
\label{prop:dT}
\begin{casos}
  \item Con la tercera ley en forma \emph{débil},
  \begin{equation}
    \od{T}{t}=\underbrace{\sa\vb{v}_{\alpha}\cdot\vb{F}_{\alpha}\ext}_{\mathcal{P}\ext}
    +\sab(\vb{v}_{\alpha}-\vb{v}_{\beta})\cdot\vb{F}_{\alpha\beta} .
    \label{eq:dT}
  \end{equation}
  \item Si además las fuerzas internas son centrales (forma \emph{fuerte}) y las
  distancias entre partículas son constantes (sólido rígido), las fuerzas
  internas no trabajan y $\dot T=\mathcal{P}\ext$.
\end{casos}
\end{proposicion}

\begin{proof}
(i) $\dot T=\sum m_{\alpha}\vb{v}_{\alpha}\cdot\dot{\vb{v}}_{\alpha}=\sum\vb{v}_{\alpha}\cdot\bigl(\vb{F}_{\alpha}\ext+\sum_{\beta\neq\alpha}\vb{F}_{\alpha\beta}\bigr)$.
Agrupando la doble suma por parejas y usando sólo $\vb{F}_{\beta\alpha}=-\vb{F}_{\alpha\beta}$,
$\vb{v}_{\alpha}\cdot\vb{F}_{\alpha\beta}+\vb{v}_{\beta}\cdot\vb{F}_{\beta\alpha}=(\vb{v}_{\alpha}-\vb{v}_{\beta})\cdot\vb{F}_{\alpha\beta}=\dot{\vb{r}}_{\alpha\beta}\cdot\vb{F}_{\alpha\beta}$.

\noindent(ii) Con la forma fuerte, $\vb{F}_{\alpha\beta}=F_{\alpha\beta}\,\uv{r}_{\alpha\beta}$, y por el \cref{lem:radial}
$\dot{\vb{r}}_{\alpha\beta}\cdot\vb{F}_{\alpha\beta}=F_{\alpha\beta}\,\dot{\vb{r}}_{\alpha\beta}\cdot\uv{r}_{\alpha\beta}=F_{\alpha\beta}\,\dot r_{\alpha\beta}$, que se anula si
$r_{\alpha\beta}=\text{const}$.
\end{proof}

\begin{observacion}[Por qué hace falta la forma fuerte]
Con la forma débil sola, $\vb{F}_{\alpha\beta}$ podría tener componente perpendicular
a $\vb{r}_{\alpha\beta}$; aunque la distancia sea fija, $\dot{\vb{r}}_{\alpha\beta}$ (el giro relativo
de la pareja) es perpendicular a $\vb{r}_{\alpha\beta}$ y esa componente trabajaría.
\end{observacion}

\begin{teorema}[Conservación de la energía de un sistema]
\label{teo:energia-sistema}
Supongamos que las fuerzas internas derivan de un potencial de pares,
\begin{equation}
  \vb{F}_{\alpha\beta}=-\od{U_{\alpha\beta}}{r_{\alpha\beta}}\,\uv{r}_{\alpha\beta},
  \qquad U\inn\defeq\sab U_{\alpha\beta}(r_{\alpha\beta}),
\end{equation}
que son centrales y conservativas. Entonces:
\begin{casos}
  \item $\od{}{t}\bigl(T+U\inn\bigr)=\mathcal{P}\ext$: sin fuerzas externas, se conserva
        la energía mecánica \emph{propia} $T+U\inn$.
  \item Si además $\vb{F}_{\alpha}\ext=-\grad_{\alpha}V\ext$, entonces
        \begin{equation}
          \resultado{E\defeq T+U\inn+V\ext=\text{const}} .
        \end{equation}
\end{casos}
\end{teorema}

\begin{notacion}[Gradiente respecto de una partícula]
$V\ext=V(\vb{r}_{1},\dots,\vb{r}_{N})$ es función de $3N$ variables. El gradiente
$\grad_{\alpha}$ deriva sólo respecto de las coordenadas de la partícula $\alpha$:
\begin{equation}
  \grad_{\alpha}V\defeq\Bigl(\pd{V}{x_{\alpha}},\ \pd{V}{y_{\alpha}},\ \pd{V}{z_{\alpha}}\Bigr),
  \qquad \vb{F}_{\alpha}\ext=-\grad_{\alpha}V\ext .
\end{equation}
Ejemplo (peso): $V\ext=\sum_{\beta}m_{\beta}g\,z_{\beta}\Rightarrow\grad_{2}V\ext=(0,0,m_{2}g)$, luego
$\vb{F}_{2}\ext=-m_{2}g\,\kk$.
\end{notacion}

\begin{proof}
Objetivo: escribir los dos sumandos de \eqref{eq:dT} como derivadas totales,
$-\dot U\inn$ y $-\dot V\ext$.

\noindent\emph{(i) Fuerzas internas}, pareja a pareja:
\begin{enumerate}
  \item \emph{Velocidad relativa}: $\vb{v}_{\alpha}-\vb{v}_{\beta}=\dot{\vb{r}}_{\alpha\beta}$ (vector).
  \item \emph{Fuerza de pares}, con $U'_{\alpha\beta}\defeq\dd U_{\alpha\beta}/\dd r_{\alpha\beta}$:
  $(\vb{v}_{\alpha}-\vb{v}_{\beta})\cdot\vb{F}_{\alpha\beta}=-U'_{\alpha\beta}\,(\dot{\vb{r}}_{\alpha\beta}\cdot\uv{r}_{\alpha\beta})$.
  \item \emph{\Cref{lem:radial}} (fuerza central: sólo trabaja la parte radial):
  $\dot{\vb{r}}_{\alpha\beta}\cdot\uv{r}_{\alpha\beta}=\dot r_{\alpha\beta}$ (escalar).
  \item \emph{Regla de la cadena} ($U_{\alpha\beta}$ depende de $t$ sólo a través de $r_{\alpha\beta}$):
  \[
    (\vb{v}_{\alpha}-\vb{v}_{\beta})\cdot\vb{F}_{\alpha\beta}=-U'_{\alpha\beta}\,\dot r_{\alpha\beta}
    =-\od{U_{\alpha\beta}}{r_{\alpha\beta}}\od{r_{\alpha\beta}}{t}=-\od{U_{\alpha\beta}}{t}.
  \]
  \item \emph{Sumando} en $\alpha<\beta$:
  \[
    \sab(\vb{v}_{\alpha}-\vb{v}_{\beta})\cdot\vb{F}_{\alpha\beta}=-\od{}{t}\sab U_{\alpha\beta}=-\dot U\inn
    \ \Longrightarrow\ \od{}{t}\bigl(T+U\inn\bigr)=\mathcal{P}\ext .
  \]
\end{enumerate}
\emph{(ii) Fuerzas externas}. Por la regla de la cadena en las $3N$ variables,
\[
  \od{V\ext}{t}=\sa\Bigl(\pd{V}{x_{\alpha}}\dot x_{\alpha}+\pd{V}{y_{\alpha}}\dot y_{\alpha}+\pd{V}{z_{\alpha}}\dot z_{\alpha}\Bigr)
  =\sa\grad_{\alpha}V\ext\cdot\vb{v}_{\alpha}=-\sa\vb{v}_{\alpha}\cdot\vb{F}_{\alpha}\ext=-\mathcal{P}\ext ,
\]
y con (i): $\dot T+\dot U\inn+\dot V\ext=0$.
\end{proof}

\begin{observacion}[Dónde entra cada hipótesis]
\begin{itemize}
  \item \emph{Centrales}: permiten el paso 3; una componente perpendicular de
        $\vb{F}_{\alpha\beta}$ trabajaría con el giro relativo (\cref{fig:radial-giro}).
  \item \emph{Conservativas}: convierten $U'_{\alpha\beta}\dot r_{\alpha\beta}$ en la derivada total del paso 4.
  \item \emph{$V\ext$ sin dependencia explícita de $t$}: si $V\ext=V\ext(\vb{r}_{1},\dots,\vb{r}_{N},t)$, la
        regla de la cadena añade $\partial V\ext/\partial t$ y $\dot E=\partial V\ext/\partial t\neq0$
        (por ejemplo, un campo externo que se enciende).
\end{itemize}
\end{observacion}

\begin{table}[ht]
\centering
\caption{Tipos de energía de un sistema de partículas.}
\label{tab:energias}
\small
\begin{tabular}{@{}l l@{}}
\toprule
\color{cian}Símbolo & \color{cian}Significado\\
\midrule
$T=\tfrac12MV^{2}+T'$ & energía cinética total (König)\\
$T'$ & energía cinética interna, respecto del CM\\
$U\inn$ & energía potencial interna (pares de partículas)\\
$V\ext$ & energía potencial de las fuerzas externas\\
$T'+U\inn$ & energía interna del sistema\\
$T+U\inn$ & energía mecánica propia: se conserva si $\vb{F}\ext=\vb{0}$\\
$T+U\inn+V\ext$ & energía mecánica total: se conserva si todo es conservativo\\
\bottomrule
\end{tabular}
\end{table}
